\documentclass[a4paper,twoside]{article}

\usepackage{morewrites}

\usepackage[svgnames,dvipsnames,x11names]{xcolor}
\usepackage{graphicx}
\usepackage[english]{babel}
\usepackage[colorlinks=true, linkcolor=NavyBlue, urlcolor=NavyBlue, citecolor=NavyBlue]{hyperref}
\usepackage[a4paper]{geometry}
\usepackage{ts-fonts}
\usepackage{tcolorbox}
\usepackage{fvextra}
\usepackage{float}
\usepackage{seqsplit}
\usepackage{ts-code}

\usepackage{lastpage}
\usepackage{refcount}
\makeatletter
\def\ps@plain{
\let\@oddhead\@empty
\let\@evenhead\@empty
\def\@oddfoot{
\hfil
\ifnum\getpagerefnumber{LastPage}>1\relax
\thepage
\fi
\hfil}
\let\@evenfoot\@oddfoot
}
\makeatother

\title{Smarty Pants}
\author{}\date{}
\begin{document}
\maketitle
This extension ports the Python-Markdown SmartyPants behavior. It swaps
plain ASCII punctuation for typographically ``smart'' equivalents---straight quotes
become curly quotes, \tscodeinline{-\allowbreak{}-\allowbreak{}} becomes an en dash, and \tscodeinline{-\allowbreak{}-\allowbreak{}-\allowbreak{}} turns into an em dash.
The substitutions happen in both HTML and \LaTeX{} output, so your prose looks
polished everywhere.
\end{document}
